\chapter{Abstract}
\onehalfspacing

Online learning platforms are now widespread in higher education, yet most of them represent student progress in the same limited way: a list of grades, or a single completion percentage for a course. Neither representation indicates which topic a student has failed to understand, nor how far that student has fallen behind the rest of the cohort. Teachers face the corresponding difficulty, since a gradebook records outcomes instead of the process that produced them, and a student who needs help is often identified only after an examination.

This report presents the design, implementation and evaluation of Synapse, a learning management system in which the course syllabus is represented as an interactive mindmap. Each topic is a node on a branching structure; a node moves from not started to in progress, and finally to completed, as the student works through the resources associated with it. Teacher and student therefore read the same representation of the course, and the position of a student within that structure is available at a glance.

The system is implemented as a Next.js application on PostgreSQL, with React Flow providing the canvas and Supabase providing authentication and file storage. The platform is multi-tenant: several institutions share one deployment while their records remain separated at the database level. A retrieval layer is provided for an AI study assistant, which answers from the documents supplied for a given subject and reports the absence of relevant material instead of producing an unsupported answer.

The contribution of the work is the combination of three mechanisms that are usually studied separately: mindmap-based progress visualisation, dual tracking of taught and learned material, and retrieval constrained to course material. The report describes the implementation of each, states the limitations of the current build, and sets out the remaining work.

\textbf{Keywords:} Learning Management System, Mindmap Visualization, EdTech, Student Progress Tracking, Gamification, Learning Analytics, Retrieval-Augmented Generation
